\textbf{Время:} 01:30

\textbf{Тезис:} C сравнивает совместную регистрацию с независимым ожиданием; независимому контролю соответствует единица.

\textbf{Визуал:} одночастичный спектр → совместное распределение → смысл C = 1.

\textbf{Текст:} Представим много событий с зарегистрированными частицами. Гистограмма одного импульса показывает, какие значения встречаются часто. Но она не отвечает, какие импульсы встречаются вместе в одном столкновении. Для этого нужно совместное распределение пары.

При независимости оно получается произведением одночастичных распределений с учётом общей нормировки. Корреляционная функция C сравнивает совместную частоту с независимым ожиданием: берём отношение двухчастичной плотности к произведению одночастичных и вводим общую нормировку N. В эксперименте используем плотности выходов частиц и пар; N выбираем так, чтобы независимому фону в контрольной области соответствовало C = 1. Значение выше единицы — выбранные импульсы встречаются вместе чаще контроля; ниже — реже. Нормированный спектр слева — лишь учебная иллюстрация одночастичного распределения.

Можно сохранить отдельные импульсы, но переставить частицы между событиями. Спектр сохранится, а парные связи изменятся. Поэтому пары дают новую информацию. Однако связь ещё не обязательно вызвана размером: существуют распады, законы сохранения и эффекты регистрации.

Пауза: прежде чем искать физический сигнал, нужно понимать значение базовой линии C = 1.

\textbf{Переход:} В какой переменной искать избыток похожих импульсов?
